\begin{scanpage}{0216}{198}
\[E_\vartheta=H_\varphi=-E\cdot\frac{e^{ikR}}{kR}\cdot F(\vartheta,\vartheta_0),\tag{32.04}\]
where
\[\begin{aligned}
F(\vartheta,\vartheta_0)={}&\frac{2i}{(\cos\vartheta+\cos\vartheta_0)\sin\vartheta\sin\vartheta_0\ln\dfrac{2i}{\gamma ka\sin\vartheta}\ln\dfrac{2i}{\gamma ka\sin\vartheta_0}}\times\\
\times\Biggl\{&\cos^2\frac\vartheta2\cos^2\frac{\vartheta_0}2\ln\frac{i}{\gamma ka\cos\dfrac\vartheta2\cos\dfrac{\vartheta_0}2}e^{-ikz_1(\cos\vartheta+\cos\vartheta_0)}\\
&-\sin^2\frac\vartheta2\sin^2\frac{\vartheta_0}2\ln\frac{i}{\gamma ka\sin\dfrac\vartheta2\sin\dfrac{\vartheta_0}2}e^{-ikz_2(\cos\vartheta+\cos\vartheta_0)}\\
&+e^{ikL}\left[\sin^2\frac\vartheta2\sin^2\frac{\vartheta_0}2\ln\frac{i}{\gamma ka\sin\dfrac{\vartheta_0}2}\psi_-\right.\\
&\qquad\left.-\cos^2\frac\vartheta2\cos^2\frac{\vartheta_0}2\ln\frac{i}{\gamma ka\cos\dfrac\vartheta2}\psi_+^0\right]e^{-ik(z_2\cos\vartheta_0+z_1\cos\vartheta)}\\
&+e^{ikL}\left[\sin^2\frac\vartheta2\sin^2\frac{\vartheta_0}2\ln\frac{i}{\gamma ka\sin\dfrac\vartheta2}\psi_-^0\right.\\
&\qquad\left.-\cos^2\frac\vartheta2\cos^2\frac{\vartheta_0}2\ln\frac{i}{\gamma ka\cos\dfrac{\vartheta_0}2}\psi_+\right]e^{-ik(z_2\cos\vartheta+z_1\cos\vartheta_0)}\\
&-\frac{(\cos\vartheta+\cos\vartheta_0)\ln\dfrac{i}{\gamma ka}}{2D}e^{2ikL}\bigl[\psi_+^0\psi e^{ikL-ikz_2\cos\vartheta_0}\\
&\qquad-\psi_-^0 e^{-ikz_1\cos\vartheta_0}\bigr]\psi_-e^{-ikz_1\cos\vartheta}\\
&-\frac{(\cos\vartheta+\cos\vartheta_0)\ln\dfrac{i}{\gamma ka}}{2D}e^{2ikL}\bigl[\psi_-^0\psi e^{ikL-ikz_1\cos\vartheta_0}\\
&\qquad-\psi_+^0 e^{-ikz_2\cos\vartheta_0}\bigr]\psi_+e^{-ikz_2\cos\vartheta}\Biggr\}.
\end{aligned}\tag{32.05}\]
in which all the functions $\psi_\pm$ and $\psi_\pm^0$ have the argument $L$. The resulting expression, despite its complexity, has a clear physical meaning. Actually, the first term in the braces corresponds to the primary edge wave radiated by the conductor's end $z=z_1$; the second term corresponds to the primary wave radiated by the conductor's end $z=z_2$. The terms included in the first pair of brackets refer to the secondary wave departing from the end $z=z_1$, and the terms in the second set of brackets refer to the secondary wave departing from the end $z=z_2$. The remaining terms describe the sum of all subsequent waves arising with multiple diffraction and have a resonance character. The resonance begins with $L=z_2-z_1\approx n\cdot\lambda/2$ $(n=1,2,3,\ldots)$ when $D\approx0$.
\end{scanpage}
